\documentclass[11pt]{article}
\usepackage[margin=1in]{geometry}
\usepackage{amsmath,amssymb,amsthm,booktabs,microtype}
\usepackage[colorlinks,linkcolor=blue!50!black,citecolor=blue!50!black]{hyperref}

\theoremstyle{definition}
\newtheorem{definition}{Definition}
\theoremstyle{plain}
\newtheorem{lemma}{Lemma}
\newtheorem{proposition}[lemma]{Proposition}
\newtheorem{conjecture}[lemma]{Conjecture}

% Every formal statement carries its epistemic status. A document that mixes
% proved, measured and guessed claims without marking which is which is worse
% than no document, because it launders one into another.
\newcommand{\PROVED}{\textnormal{\footnotesize\textsf{[PROVED]}}}
\newcommand{\MEASURED}{\textnormal{\footnotesize\textsf{[MEASURED]}}}
\newcommand{\OPEN}{\textnormal{\footnotesize\textsf{[OPEN]}}}

\newcommand{\wta}{\mathrm{top}_k}
\newcommand{\E}{\mathbb{E}}
\newcommand{\Prob}{\mathbb{P}}

\title{Statistical Mechanics of Assembly Calculus:\\
       Order Parameters, Control Parameters, and a Measured Critical Gain}
\author{Assemblies project --- working document}
\date{2026-07-29}

\begin{document}
\maketitle

\begin{abstract}
We set up the assembly-calculus substrate as a statistical-mechanical system:
microstates (the connectome and the firing pattern), macrostates (overlap,
retrieval accuracy, margin), and control parameters (load $\alpha$, gain $g$,
afferent number $\mu$, depth $D$). Within this frame, a collapse previously
described as a ``depth ceiling'' is identified as a transition in a single
control parameter, the \emph{gain} $g=(1+\beta)^T$, with a measured critical
value $g_c \approx 1.9$ that is invariant to how $\beta$ and $T$ factorise it
and invariant to connection density $p$. Statements below are marked
\PROVED, \MEASURED\ or \OPEN; nothing is asserted as proved that has only been
observed.
\end{abstract}

\section{The substrate}

\begin{definition}[Area and assembly]
An \emph{area} $A$ consists of $n$ neurons with a firing cap $k \ll n$. A
\emph{state} of $A$ is a set $x \subseteq [n]$ with $|x| = k$. There are
$\binom{n}{k}$ such states; we call this the \emph{configuration space} of $A$.
\end{definition}

\begin{definition}[Connectome and microstate]
A \emph{fiber} from area $A$ to area $B$ carries weights $W \in
\mathbb{R}^{n_B \times n_A}$. At initialisation $W_{ij} \sim
\mathrm{Bernoulli}(p)$ independently. The \emph{microstate} of the system is
the pair $(\{W^{(f)}\}_f, \{x_A\}_A)$: every synaptic weight together with the
current firing set of every area.
\end{definition}

\begin{definition}[Dynamics]
In one round, neuron $i \in B$ receives
\(
  I_i = \sum_{A \to B} \sum_{j \in x_A} W^{(A\to B)}_{ij},
\)
and the new state is $x_B' = \wta(I)$, the $k$ largest coordinates (ties broken
by index). Plasticity is multiplicative and Hebbian: for $j \in x_A$ and $i \in
x_B'$,
\[
  W_{ij} \;\longleftarrow\; \min\!\big(W_{ij}\,(1+\beta),\; w_{\max}\big).
\]
\end{definition}

\begin{definition}[Gain]\label{def:gain}
If a source assembly is presented for $T$ consecutive rounds and the target
assembly is stable across them, every fiber from source to target is
potentiated $T$ times. The resulting multiplicative advantage
\[
  \boxed{\,g \;=\; (1+\beta)^T\,}
\]
is the \emph{gain}. It is the natural coupling constant of the system: $\beta$
and $T$ enter the dynamics only through this product over a stable
presentation.
\end{definition}

\section{Macrostates: order parameters}

\begin{definition}[Overlap]
For assemblies $a,b$ in the same area, $q_{ab} = |a \cap b| / k$. For a family
$\{a_m\}_{m=1}^{M}$ the \emph{spread} is the mean pairwise overlap
$\bar q = \binom{M}{2}^{-1}\sum_{m<m'} q_{a_m a_{m'}}$.
\end{definition}

\begin{lemma}[Overlap floor]\label{lem:floor} \PROVED
If $a,b$ are independent uniformly random $k$-subsets of $[n]$ then
$\E[q_{ab}] = k/n$.
\end{lemma}
\begin{proof}
$|a \cap b| = \sum_{i=1}^{n} \mathbf{1}[i \in a]\mathbf{1}[i \in b]$. The two
indicators are independent with means $k/n$, so $\E|a\cap b| = n(k/n)^2 =
k^2/n$, and dividing by $k$ gives $k/n$.
\end{proof}

The floor $q_0 = k/n$ is size dependent, so raw overlap cannot be compared
across system sizes. We therefore use the normalised order parameter
\begin{equation}\label{eq:Q}
  Q \;=\; \frac{\bar q - q_0}{1 - q_0} \in [0,1],
\end{equation}
which vanishes for mutually independent assemblies at any $n$ and equals $1$
under total collapse. $Q$ plays the role of the Edwards--Anderson overlap in
spin glasses, and of the retrieval overlap in the Hopfield model.

Two further macrostates are measured. \emph{Retrieval accuracy} $R$ is the
probability that a re-presented item's read-out has its largest overlap with
the correct stored assembly. \emph{Margin} $\rho$ is the ratio of best to
second-best overlap; it is the more sensitive statistic, moving well before $R$
departs from $1$.

\section{Control parameters}

\begin{align}
  \alpha &= \frac{Mk}{n} && \text{\emph{load}: fraction of the substrate committed}\\
  g &= (1+\beta)^T && \text{\emph{gain}: Definition~\ref{def:gain}}\\
  \mu &= kp && \text{\emph{afferent number}: mean inputs from one assembly}\\
  D && &\text{\emph{depth}: levels of composition}
\end{align}

The load $\alpha$ is the exact analogue of the Hopfield storage ratio. The
afferent number $\mu$ has no Hopfield counterpart because that model is fully
connected; here it is decisive, and its role is the subject of
Proposition~\ref{prop:mu}.

\section{Measured phase structure}

At $\alpha = 0.8$, $\mu = 2.5$, $D = 5$, three regimes are distinguished by the
\emph{pair} $(Q, R)$ --- neither alone separates them.

\begin{center}
\begin{tabular}{llll}
\toprule
Phase & $Q$ & $R$ & Mechanism \\
\midrule
Starved   & $\approx 0$ & $<1$ & assemblies distinct but drive too weak to retrieve\\
Retrieval & $\approx 0$ & $\approx 1$ & the working phase \\
Collapsed & $\to 1$ & $\to 0$ & assemblies merge on contact \\
\bottomrule
\end{tabular}
\end{center}

That the starved and collapsed phases are \emph{both} low-$R$ while differing
in $Q$ is the practically important fact: they require opposite corrections
(raise $g$ versus lower $g$), and a diagnostic keyed on $R$ alone cannot tell
them apart.

\begin{conjecture}[Critical gain]\label{conj:gc} \MEASURED
For $T \ge 2$ and at fixed $(n,\alpha,\mu,D)$ there is a critical gain $g_c$
such that composition to depth $5$ succeeds for $g < g_c$ and fails for
$g > g_c$, and $g_c$ depends on $\beta$ and $T$ only through $g$. At $n=4000$,
$\alpha=0.8$, $\mu=2.5$, $D=5$ the boundary lies at $g_c \approx 1.9$.
\end{conjecture}

\paragraph{$g_c$ is not a constant of the substrate.} The size qualification
above is not pedantry. A fine scan at matched load gives, by the $R=1/2$
crossing, $g_c \approx 1.71$ at $n=2000$ against $\approx 1.9$ at $n=4000$: the
critical gain \emph{rises with system size}. This is consistent in direction
with the separately measured threshold for recurrence persistence, which also
rises with $n$. Any statement of the form ``$g_c \approx 2$'' is therefore a
statement about one system size, and the object of interest is the function
$g_c(n,\alpha,\mu,D)$ --- which is what Conjecture~\ref{conj:fss} is really
about, and what makes the $\sqrt{2\ln(n/k)}$ term in
Conjecture~\ref{conj:why} the right shape to be looking for.

Evidence ($n=4000$, $\alpha=0.8$, $\mu=2.5$, $D=5$, 2 seeds): $\beta$ ranges
over $0.18$--$0.45$ across the boundary while $g$ does not move.

\begin{center}
\begin{tabular}{ccccc}
\toprule
$T$ & $\beta$ & $g=(1+\beta)^T$ & $R$ & $\bar q$ \\
\midrule
3 & 0.18 & 1.643 & 1.000 & 0.0157 \\
2 & 0.30 & 1.690 & 1.000 & 0.0180 \\
3 & 0.22 & 1.816 & 0.883 & 0.0254 \\
2 & 0.35 & 1.823 & 0.938 & 0.0247 \\
2 & 0.40 & 1.960 & 0.164 & 0.0812 \\
3 & 0.26 & 2.000 & 0.117 & 0.2359 \\
2 & 0.45 & 2.103 & 0.070 & 0.2691 \\
3 & 0.32 & 2.300 & 0.047 & 0.6170 \\
\bottomrule
\end{tabular}
\end{center}

$T=1$ lies systematically below this curve (at $g=1.800$ it gives $R=0.625$
where $T\in\{2,3\}$ give $0.938$ and $0.883$), so the collapse onto $g$ is
asserted only for $T \ge 2$. A single round does not let the target assembly
converge before plasticity is applied, which is a separate effect and is not
described by $g$.

\begin{proposition}[Density sets the level, not the location]\label{prop:mu} \MEASURED
Doubling $p$ from $0.05$ to $0.10$ (i.e.\ $\mu$ from $2.5$ to $5$) raises the
depth-5 margin from $4.65$ to $8.11$ and $R$ from $0.938$ to $1.000$ at
$\beta=0.35$, while leaving the failure at $\beta=0.40$ unchanged
($R = 0.164 \to 0.172$). Moreover the \emph{sign} of $\partial \rho/\partial
\beta$ depends on $\mu$: for $\mu \lesssim 1$ the margin has an interior
maximum in $\beta$, and for $\mu \gtrsim 2.5$ it is increasing over the tested
range.
\end{proposition}

Interpretation: with few afferents, potentiation concentrates drive on a
handful of neurons that then win every subsequent competition, so large $\beta$
manufactures collapse; with many afferents the same potentiation is spread
across a rich pattern and selectivity survives. Density therefore controls the
\emph{amplitude} of the signal but not the \emph{stability boundary}.

\section{What remains to be derived}

\subsection{An empirical law for \texorpdfstring{$g_c$}{g\_c}}

Untangling $n$ from $M$ --- the main scan holds $\alpha$ fixed and so moves
both at once --- gives three cells that separate the two variables:

\begin{center}
\begin{tabular}{ccccl}
\toprule
$n$ & $M$ & $\alpha$ & $g_c$ & \\
\midrule
2000 & 32 & 0.8 & 1.715 & \\
4000 & 32 & 0.4 & 2.260 & $n$ doubled at fixed $M$: $\Delta = +0.545$\\
4000 & 64 & 0.8 & 1.917 & $M$ doubled at fixed $n$: $\Delta = -0.343$\\
2000 & 64 & 1.6 & $<1.50$ & below the scanned range\\
\bottomrule
\end{tabular}
\end{center}

Both variables move $g_c$, and in opposite directions, so neither the
fresh-pool picture ($n$ only) nor the stored-set picture ($M$ only) is
complete on its own. Reading the two increments as logarithmic gives
$\partial g_c/\partial \ln n \approx 0.79$ and
$\partial g_c/\partial \ln M \approx -0.50$, and substituting $M = \alpha n/k$
collapses these to

\begin{equation}\label{eq:gclaw}
  g_c \;\approx\; C \;+\; 0.29\,\ln n \;-\; 0.50\,\ln \alpha .
\end{equation}

\begin{conjecture}[Empirical form of $g_c$]\label{conj:gclaw} \MEASURED
$g_c$ increases logarithmically with system size at fixed load and decreases
logarithmically with load, as in \eqref{eq:gclaw}.
\end{conjecture}

Three points determine three coefficients, so \eqref{eq:gclaw} is not yet
tested --- it is a fit. Its content is the \emph{prediction} it makes at a
fourth cell. For $n=8000$, $M=128$ (again $\alpha=0.8$, both doubling once
more) it requires
\[
  g_c(8000,128) \;\approx\; 1.917 + 0.29\ln 2 \;\approx\; \boxed{2.12}.
\]
This is registered before that cell finished measuring. A result near $2.12$
supports the form; convergence towards a constant instead would mean $g_c$ has
a thermodynamic limit and the logarithm is a small-$n$ artifact; continued
growth at a different rate would falsify the coefficient.

The sign structure is worth stating plainly, because it is the practical
content: \textbf{bigger substrates tolerate more gain, and heavier loads
tolerate less.} A gain that is safe on a small vocabulary can therefore become
unsafe purely by adding items --- with no change to any plasticity parameter.

\begin{conjecture}[Finite-size scaling]\label{conj:fss} \OPEN
Near $g_c$ the order parameter obeys
\[
  Q(g,n) \;=\; n^{-\gamma/\nu}\, F\!\big((g-g_c)\, n^{1/\nu}\big)
\]
for a scaling function $F$ and exponents $\gamma,\nu$; equivalently, curves of
$Q$ against $g$ at different $n$ cross at $g_c$ and collapse when replotted
against $(g-g_c)n^{1/\nu}$.
\end{conjecture}

This is the statement that distinguishes a genuine continuous transition from a
mere threshold, and it is what \texttt{critical\_point\_scan.py} is measuring.
A threshold set by a fixed arithmetic condition would show no $n$ dependence at
all; a first-order transition would show hysteresis instead of collapse.

\begin{conjecture}[Why $g_c \approx 2$]\label{conj:why} \OPEN
The observed $g_c$ lies close to $2$, the point at which one stable
presentation doubles a fiber's weight. A derivation should compare the
potentiated drive of a stored assembly's members against the upper order
statistics of the fresh candidate pool: with afferent counts
$\mathrm{Binomial}(k,p)$, the $k$-th largest of $n$ candidates sits
approximately $\sqrt{2\ln(n/k)}$ standard deviations above the mean, giving a
condition of the form $g > 1 + c\sqrt{(1-p)/(kp)}\,\sqrt{2\ln(n/k)}$ for a
veteran to displace fresh candidates.
\end{conjecture}

We flag a tension that must be resolved rather than glossed: the right-hand
side of that expression depends on $p$, whereas $g_c$ was measured to be
\emph{independent} of $p$ (Proposition~\ref{prop:mu}). Either the relevant
comparison is not against fresh candidates but against \emph{other stored}
assemblies --- whose drive is potentiated by the same $g$, cancelling the
density dependence --- or the argument is wrong. The first alternative is the
more likely and is the natural next derivation.

\begin{proposition}[Depth is information-limited]\label{prop:dpi} \OPEN
Model level $L$ of a chain as a channel $X \to A_1 \to \cdots \to A_D$ where
$X$ is the item identity, uniform on $M$ values. By the data-processing
inequality $I(X; A_1) \ge \cdots \ge I(X; A_D)$, and exact retrieval at depth
$D$ requires $I(X;A_D) = \log M$. Hence if each level loses at least $\delta$
bits, the maximum reliably composable depth is at most
$(\log M - \log M_{\min})/\delta$.
\end{proposition}

This turns depth decay into a measurable quantity --- $\delta$ is estimable
from the confusion matrix at each level --- and gives a bound that does not
depend on the details of $\wta$. Making it quantitative requires estimating
$I(X;A_L)$ from the empirical read-out distribution, which the scan's per-level
records already support.

\section{Operating point: is the substrate quasicritical?}

A measured critical point invites the question the critical-brain literature
asks of cortex: does the system work best \emph{at} criticality, or near it?
Here the question is sharp because $g$ is an explicit parameter rather than
something the system must self-organise towards.

\paragraph{A scope caveat, stated plainly.} This substrate is not a spiking
network. It has no membrane dynamics, no continuous time, and no avalanches;
one ``round'' is a synchronous $k$-WTA update. So the classical signatures of
critical brain dynamics --- avalanche size and duration distributions with
exponents near $-3/2$ and $-2$, and the associated crackling-noise relation ---
are \emph{not} measurable here, and any claim of continuity with those results
would be an equivocation. What transfers is the weaker and still useful
statement: a control parameter with a critical value, an order parameter that
changes character across it, and an operating point defined relative to it.

\begin{proposition}[The optimum lies strictly below $g_c$]\label{prop:sub} \MEASURED
At depth $5$, margin is not maximised as $g \to g_c^-$. At $T=2$, $g=1.690$
gives margin $6.53$ while $g=1.823$ gives $4.65$; at $T=3$, $g=1.643$ gives
$9.65$ while $g=1.816$ gives $4.85$. Performance degrades on approach to the
transition from below, before the transition itself.
\end{proposition}

This is the opposite of ``tune to the critical point''. It suggests the useful
regime is \emph{subcritical with headroom}, and it has a natural reading: the
approach to $g_c$ is where overlap begins to accumulate, and at depth that
accumulation compounds multiplicatively along the chain.

\begin{conjecture}[Depth lowers the optimal gain]\label{conj:depth-gain} \OPEN
The gain maximising margin decreases with composition depth $D$. At $D=3$ the
margin is still increasing at $\beta=0.35$ ($4.70 \to 7.89 \to 8.97$ for
$\beta = 0.10, 0.20, 0.35$); at $D=5$ an interior maximum appears near
$\beta \approx 0.30$.
\end{conjecture}

The evidence for Conjecture~\ref{conj:depth-gain} is currently assembled across
runs with differing seeds and is the weakest claim in this document; it is
stated because it is directly testable on the scan data, not because it is
established.

\paragraph{Consequence: a per-area gain schedule.} If the optimal gain falls
with depth, then a single global $\beta$ is provably suboptimal for a
hierarchy, and areas deeper in the composition chain should carry lower gain.
Since $g=(1+\beta)^T$, this can be realised either by lowering $\beta$ per area
or by shortening $T$ per level --- the two are interchangeable to the extent
Conjecture~\ref{conj:gc} holds, which is itself a testable prediction about a
design choice. A time-varying schedule (high gain early to recruit, decaying
later to protect what is stored) is the same statement along the training axis
rather than the spatial one, and is the natural analogue of a
neuromodulatory gating of plasticity.

\section{Ensembles and entropy}

The configuration entropy of one area is
\(
  S \;=\; \ln\binom{n}{k} \;\approx\; k\ln(n/k) + k
\)
nats, an upper bound on what an area can represent. Capacity questions then
become counting arguments: the $M$ stored assemblies must be mutually
distinguishable under the read-out, and $\alpha$ measures how much of the
configuration space is committed.

A canonical ensemble can be introduced by softening $\wta$: assign to each
state $x$ the energy $E(x) = -\sum_{i \in x} I_i$ and the Gibbs measure
$\Prob(x) \propto e^{-E(x)/\tau}$ over $k$-subsets, with partition function
$Z(\tau) = \sum_{|x|=k} e^{-E(x)/\tau}$. Deterministic $k$-WTA is the
$\tau \to 0$ limit. This is the natural route to a free energy for the
substrate, and it is worth noting that $\tau$ is not currently a parameter of
the implementation --- introducing it would be a modelling choice, not a
measurement, and should be kept clearly separate from the results above.

\end{document}
